\documentclass[11pt]{article}

\usepackage[margin=1in]{geometry}
\usepackage{amsmath,amssymb,amsfonts}
\usepackage{booktabs}
\usepackage{longtable}
\usepackage{array}
\usepackage{hyperref}
\usepackage{enumitem}

\hypersetup{
  colorlinks=true,
  linkcolor=blue,
  urlcolor=blue,
  citecolor=blue
}

\title{\textbf{Metric Mathematical Derivations for EquiMed-DSS}\\
\large Formulae Matched to the EquiMed-DSS Library Implementation}
\author{John Weirstrass Muteba Mwamba\\
\small School of Health Policy and Management, York University, Toronto, Canada\\
\small Developer and Modeler of EquiMed-DSS\\
\small \href{mailto:jwmm@yorku.ca}{jwmm@yorku.ca}}
\date{EquiMed-DSS version 1.5.1}

\newcommand{\ind}{\mathbb{I}}
\newcommand{\E}{\mathbb{E}}
\newcommand{\Var}{\mathrm{Var}}
\newcommand{\Cov}{\mathrm{Cov}}
\newcommand{\KL}{D_{\mathrm{KL}}}
\newcommand{\R}{\mathbb{R}}
\newcommand{\norm}[1]{\left\lVert #1 \right\rVert}

\begin{document}
\maketitle

\section*{Purpose and Verification Scope}

This technical document derives the mathematical definitions used by the
EquiMed-DSS library, version 1.5.1. The formulae below are aligned with the
local implementation in the package source code, especially the modules
\texttt{domain1}, \texttt{domain2}, \texttt{domain3}, \texttt{domain4},
\texttt{domain5}, \texttt{geographic}, \texttt{appendix}, and
\texttt{statistics}. Where an earlier derivation document used a more general
or theoretical expression, the implementation-specific expression is reported
here as the authoritative formula.

The library documentation lists 37 named metrics. This document gives 39
derivations by adding two implemented statistical estimands used with the
metric suite: hierarchical variance partitioning and mediation/proportion
mediated. These two estimands are implemented in \texttt{equimed\_dss.statistics}
and are included because they are used to interpret fairness and governance
results in the EquiMed-DSS framework.

Each metric is numbered in the overview table and in its subsection title. Some
implemented classes return companion quantities, for example HER with Bias-Gini
or CPS with CFU. These companion quantities use the parent metric number with a
letter suffix. Every displayed equation is written in a numbered LaTeX
\texttt{equation} or \texttt{align} environment so it can be referenced directly
after compilation.

\section*{Notation}

Let $i=1,\ldots,n$ index observations, $g \in \mathcal{G}$ index demographic or
intersectional groups, $s \in \mathcal{S}$ index healthcare-system strata, and
$r \in \mathcal{R}$ index geographic regions. Let $\hat{Y}_i$ denote a model
prediction, $Y_i$ the corresponding reference outcome, $C_i$ a confidence score,
and $Z_i=\ind(\hat{Y}_i=Y_i)$ a correctness indicator. For text outputs, let
$R_i$ denote a response, $D(R_i)$ the set of diagnoses mentioned in the response,
and $D^\star$ a reference differential diagnosis set. Unless otherwise stated,
all means are empirical means over the supplied input arrays.

\begin{longtable}{p{0.08\linewidth}p{0.32\linewidth}p{0.18\linewidth}p{0.32\linewidth}}
\toprule
No. & Metric & Abbreviation & Canonical implementation \\
\midrule
\endhead
1 & Inter-rater reliability & ICC(2,1) & \texttt{domain1.icc} \\
2 & Embedding consistency score & ECS & \texttt{domain1.ecs} \\
3 & Decision flip rate & DFR & \texttt{domain1.dfr} \\
4 & Hierarchical equity ratio and Bias-Gini & HER & \texttt{domain2.her} \\
5 & Harm-adjusted fairness gap & HAFG & \texttt{domain2.hafg} \\
6 & Ethical risk index and safety violation rate & ERI & \texttt{domain2.eri} \\
7 & Intersectional bias score & IBS & \texttt{domain2.ibs} \\
8 & Temporal fairness drift & TFD & \texttt{domain3.tfd} \\
9 & Audit traceability score & ATS & \texttt{domain3.ats} \\
10 & Governance compliance index & GCI & \texttt{domain3.gci} \\
11 & Semantic parity gap & SPG & \texttt{domain4.spg} \\
12 & Clinical hallucination rate & CHR & \texttt{domain4.chr} \\
13 & Instructional vulnerability index & IVI & \texttt{domain4.ivi} \\
14 & Geographic representation index and geographic bias & GRI & \texttt{domain4.gri} \\
15 & Intersectional calibration error & ICE & \texttt{domain5.calibration} \\
16 & Weighted clinical harm-adjusted fairness gap & wHAFG & \texttt{domain5.harm} \\
17 & Counterfactual parity score and counterfactual unfairness & CPS & \texttt{domain5.counterfactual} \\
18 & Semantic robustness parity index & SRPI & \texttt{domain5.counterfactual} \\
19 & Lexical diversity disparity index & LDDI & \texttt{domain5.text} \\
20 & Recommendation entropy gap & REG & \texttt{domain5.text} \\
21 & Clinical information density ratio & CIDR & \texttt{domain5.text} \\
22 & Diagnostic completeness index & DCI & \texttt{domain5.text} \\
23 & Uncertainty quantification gap & UQG & \texttt{domain5.text} \\
24 & Geographic representation bias index & GRBI & \texttt{domain5.geographic\_bias} \\
25 & Healthcare system stratified fairness & HSSF & \texttt{domain5.system} \\
26 & Intersectional Shapley fairness value & ISFV & \texttt{domain5.shapley} \\
27 & Burden-evidence mismatch index & BEMI & \texttt{geographic.burden\_evidence} \\
28 & Geographic concentration of coverage & GCC & \texttt{geographic.concentration} \\
29 & Bootstrap confidence interval & BCI & \texttt{appendix.advanced\_metrics} \\
30 & Statistical power analysis & SPA & \texttt{appendix.advanced\_metrics} \\
31 & Bias concentration index & BCI-bias & \texttt{appendix.advanced\_metrics} \\
32 & Mutual information content & MIC & \texttt{appendix.advanced\_metrics} \\
33 & Jensen-Shannon divergence & JSD & \texttt{appendix.advanced\_metrics} \\
34 & Wasserstein distance & WD & \texttt{appendix.advanced\_metrics} \\
35 & Network modularity & NM & \texttt{appendix.advanced\_metrics} \\
36 & Transparency score & TS & \texttt{appendix.advanced\_metrics} \\
37 & Robustness certification score & RCS & \texttt{appendix.advanced\_metrics} \\
38 & Hierarchical variance partitioning & HLM/VPC & \texttt{statistics.hierarchical} \\
39 & Causal mediation and proportion mediated & PM & \texttt{statistics.mediation} \\
\bottomrule
\end{longtable}

\section{Domain 1: Reliability and Robustness}

\subsection{Metric 1: Inter-rater Reliability, ICC(2,1)}

Let $X_{ij}$ be the score assigned to item $i$ by judge $j$, with $n$ items and
$k$ judges. Define the grand mean $\bar{X}_{..}$, item means $\bar{X}_{i.}$,
and judge means $\bar{X}_{.j}$. This follows the Shrout-Fleiss ICC family and
the Bland-Altman agreement convention \cite{shrout1979,blandaltman1986}. The
implementation computes
\begin{align}
SS_{\mathrm{items}} &= k \sum_{i=1}^{n}(\bar{X}_{i.}-\bar{X}_{..})^2,\\
SS_{\mathrm{judges}} &= n \sum_{j=1}^{k}(\bar{X}_{.j}-\bar{X}_{..})^2,\\
SS_{\mathrm{error}} &= \sum_{i=1}^{n}\sum_{j=1}^{k}(X_{ij}-\bar{X}_{..})^2
 - SS_{\mathrm{items}} - SS_{\mathrm{judges}}.
\end{align}
The mean squares are
\begin{align}
MS_R &= \frac{SS_{\mathrm{items}}}{n-1},&
MS_C &= \frac{SS_{\mathrm{judges}}}{k-1},&
MS_E &= \frac{SS_{\mathrm{error}}}{(n-1)(k-1)}.
\end{align}
EquiMed-DSS implements the two-way random-effects, single-measure, absolute
agreement intraclass correlation coefficient as
\begin{equation}
ICC(2,1)=
\frac{MS_R-MS_E}
{MS_R+(k-1)MS_E+\frac{k}{n}(MS_C-MS_E)}.
\end{equation}
The same class also computes Bland-Altman pairwise agreement. For two judges
$a$ and $b$,
\begin{align}
d_i &= X_{ia}-X_{ib},\\
\bar{d} &= \frac{1}{n}\sum_{i=1}^{n}d_i,\\
s_d &= \sqrt{\frac{1}{n-1}\sum_{i=1}^{n}(d_i-\bar{d})^2},\\
LOA_{\mathrm{lower}}, LOA_{\mathrm{upper}} &= \bar{d}\pm 1.96s_d.
\end{align}

\subsection{Metric 2: Embedding Consistency Score}

For paired original and perturbed embedding vectors $e_i$ and $\tilde e_i$, the
library calculates cosine distance, not cosine similarity:
\begin{equation}
ECS_i = 1-\frac{e_i^\top \tilde e_i}{\norm{e_i}_2\norm{\tilde e_i}_2}.
\end{equation}
If either vector has zero norm, the implementation sets the cosine similarity
to zero, so $ECS_i=1$. The reported summary statistics are
\begin{align}
\overline{ECS} &= \frac{1}{n}\sum_{i=1}^{n}ECS_i,\\
SD(ECS) &= \sqrt{\frac{1}{n}\sum_{i=1}^{n}(ECS_i-\overline{ECS})^2},\\
\widetilde{ECS} &= \mathrm{median}(ECS_1,\ldots,ECS_n).
\end{align}
Lower values indicate greater embedding stability.

\subsection{Metric 3: Decision Flip Rate}

Let $A_i$ be the model decision on the original input and $B_i$ the decision on
the paired counterfactual or perturbed input. EquiMed-DSS defines
\begin{equation}
DFR = \frac{1}{n}\sum_{i=1}^{n}\ind(A_i\ne B_i).
\end{equation}
If $x=\sum_i\ind(A_i\ne B_i)$ and $\hat p=x/n$, the Wilson interval returned by
the library is
\begin{align}
d &= 1+\frac{z^2}{n},\\
c &= \frac{\hat p+\frac{z^2}{2n}}{d},\\
h &= \frac{z\sqrt{\frac{\hat p(1-\hat p)}{n}+\frac{z^2}{4n^2}}}{d},\\
CI_{95\%} &= [\max(0,c-h),\min(1,c+h)],
\end{align}
with $z=1.96$, following Wilson's score interval for binomial proportions
\cite{wilson1927}.

\section{Domain 2: Fairness, Equity, and Ethics}

\subsection{Metric 4: Hierarchical Equity Ratio}

Let $q_g$ be a group-specific performance score and $q_0$ the reference-group
score. The implementation calculates
\begin{equation}
HER_g = \frac{q_g}{q_0},
\end{equation}
with $HER_g=0$ if $q_0=0$. Values in $[0.8,1.25]$ are labelled equitable by the
implemented four-fifths-rule convention \cite{uniformguidelines1978}.

\subsection{Metric 4a: Bias-Gini Dispersion}

For group scores $q_1,\ldots,q_K$ with mean $\bar q$, the implemented dispersion
index is the standard Gini coefficient \cite{gini1912}:
\begin{equation}
G_{\mathrm{bias}}=
\frac{\sum_{i=1}^{K}\sum_{j=1}^{K}|q_i-q_j|}
{2K^2\bar q}.
\end{equation}
If the score list is empty or $\bar q=0$, the function returns zero.

\subsection{Metric 5: Harm-adjusted Fairness Gap}

For two groups, let $FN_g$ and $FP_g$ be false-negative and false-positive
counts. With default costs $c_{FN}=10$ and $c_{FP}=3$, group harm is
\begin{equation}
H_g = c_{FN}FN_g+c_{FP}FP_g.
\end{equation}
EquiMed-DSS reports the absolute gap
\begin{equation}
\Delta_H = |H_1-H_2|
\end{equation}
and the normalized harm-adjusted fairness gap
\begin{equation}
HAFG=\frac{|H_1-H_2|}{\max(H_1,H_2)}.
\end{equation}
If both group harms are zero, the denominator is zero and the implemented value
is $0$.

\subsection{Metric 6: Ethical Risk Index}

Let $v=1,\ldots,V$ index detected ethical or safety violations and let $s_v$ be
their severity scores. For $N$ total model outputs,
\begin{equation}
ERI = \frac{\sum_{v=1}^{V}s_v}{N}.
\end{equation}
If $N=0$, the function returns zero.

\subsection{Metric 6a: Safety Violation Rate}

In the same implementation as ERI, the safety violation rate is
\begin{equation}
SVR = 1000\frac{V}{N},
\end{equation}
that is, violations per 1000 model outputs.

\subsection{Metric 7: Intersectional Bias Score}

Let $u_g\in\R^d$ be a vector of metrics for subgroup $g$. The implementation
calculates Euclidean distances
\begin{equation}
d(g,h)=\norm{u_g-u_h}_2
\end{equation}
and converts them to similarities by inverse distance:
\begin{equation}
S(g,h)=\frac{1}{1+d(g,h)}.
\end{equation}
The outlier subgroup is the subgroup with the largest mean distance over the
full distance row, including the zero self-distance used by the implementation:
\begin{equation}
g^\star=\arg\max_g \frac{1}{K}\sum_{h\in\mathcal{G}} d(g,h).
\end{equation}
The same class also computes simplified eta-squared-style interaction effects.
For a categorical attribute $A$, with grand mean $\bar Y$ and group means
$\bar Y_a$, the main-effect proxy is
\begin{equation}
\eta_A^2 =
\frac{\sum_a n_a(\bar Y_a-\bar Y)^2}
{\sum_i(Y_i-\bar Y)^2}.
\end{equation}
For the race-by-gender interaction, the implementation subtracts the race and
gender main-effect proxies from the combined race-gender proxy.

\section{Domain 3: Governance and Transparency}

\subsection{Metric 8: Temporal Fairness Drift}

For a time series of fairness metric values $m_1,\ldots,m_T$, EquiMed-DSS
computes
\begin{align}
\bar m &= \frac{1}{T}\sum_{t=1}^{T}m_t,\\
s_m &= \sqrt{\frac{1}{T-1}\sum_{t=1}^{T}(m_t-\bar m)^2}.
\end{align}
The three-sigma control limits are
\begin{align}
UCL &= \bar m+3s_m,\\
LCL &= \bar m-3s_m.
\end{align}
A drift point is flagged when $m_t>UCL$ or $m_t<LCL$, following the
Shewhart-style control-chart logic used in statistical process control
\cite{shewhart1931}.

\subsection{Metric 9: Audit Traceability Score}

Let $x$ be the number of traceable decisions among $n$ audited decisions. The
score is
\begin{equation}
ATS=\frac{x}{n}.
\end{equation}
The implementation returns a Wilson-style shrinkage interval using
\begin{align}
z &= 1.96,\\
\tilde p &= \frac{x+z^2/2}{n+z^2},\\
SE_{\tilde p} &= \sqrt{\frac{\tilde p(1-\tilde p)}{n+z^2}},\\
CI_{95\%} &= [\max(0,\tilde p-zSE_{\tilde p}),\min(1,\tilde p+zSE_{\tilde p})].
\end{align}
The interval uses Wilson score shrinkage for a binomial proportion
\cite{wilson1927}.
The score is labelled compliant at $ATS\ge 0.95$.

\subsection{Metric 10: Governance Compliance Index}

Let $M$ be the number of mandated policies and $E$ the number evaluated as
enforced. The implementation defines
\begin{equation}
GCI=\frac{E}{M}.
\end{equation}
When no policies are provided, the returned value is zero.

\section{Domain 4: Representation and Robustness}

\subsection{Metric 11: Semantic Parity Gap}

Let $P=\{p_i\}_{i=1}^{n_p}$ be embeddings for privileged-group prompts and
$M=\{m_j\}_{j=1}^{n_m}$ embeddings for matched marginalized-group prompts.
Their centroids are
\begin{align}
\bar p &= \frac{1}{n_p}\sum_{i=1}^{n_p}p_i,&
\bar m &= \frac{1}{n_m}\sum_{j=1}^{n_m}m_j.
\end{align}
The Euclidean SPG reported by the implementation is
\begin{equation}
SPG_{\mathrm{Euc}}=\norm{\bar p-\bar m}_2.
\end{equation}
The cosine variant is
\begin{equation}
SPG_{\mathrm{cos}}=1-\frac{\bar p^\top \bar m}{\norm{\bar p}_2\norm{\bar m}_2},
\end{equation}
with value zero if the denominator is zero.

\subsection{Metric 12: Clinical Hallucination Rate}

Let $c=1,\ldots,C$ index extracted clinical claims and $S(c,K)\in[0,1]$ be a
precomputed support score against retrieved context $K$. With entailment
threshold $\tau$, the implemented hallucination rate is
\begin{equation}
CHR=\frac{1}{C}\sum_{c=1}^{C}\ind\{S(c,K)<\tau\}.
\end{equation}
With optional severity weights $w_c$,
\begin{equation}
CHR_w=
\frac{\sum_{c=1}^{C}w_c\ind\{S(c,K)<\tau\}}
{\sum_{c=1}^{C}w_c}.
\end{equation}
If no weights are supplied, the implementation sets $CHR_w=CHR$.

\subsection{Metric 13: Instructional Vulnerability Index}

Let $A_i$ be the neutral-output decision and $B_i$ the paired biased-instruction
decision for the same case. The flip component is
\begin{equation}
IVI=\frac{1}{n}\sum_{i=1}^{n}\ind(A_i\ne B_i).
\end{equation}
If the outputs can be coerced to numeric values, the implementation also returns
the directional effect
\begin{equation}
IVI_{\mathrm{effect}}=\frac{1}{n}\sum_{i=1}^{n}B_i-\frac{1}{n}\sum_{i=1}^{n}A_i.
\end{equation}
The paired-counterfactual framing is conceptually related to counterfactual
fairness, although IVI targets prompt framing rather than protected-attribute
interventions \cite{kusner2017}.

\subsection{Metric 14: Geographic Representation Index}

Let $L$ be the set of unique locations represented in a corpus and let
$W\subseteq L$ be those counted as Western or high-income. The implemented index
is set-based:
\begin{equation}
GRI=\frac{|L|-|W|}{|L|}.
\end{equation}
Duplicates do not change the score.

\subsection{Metric 14a: Geographic Bias Correlation}

For paired values $(x_i,y_i)$, where $x_i$ is a per-query GRI value and $y_i$ is
the corresponding non-Western error rate, the library returns either Pearson or
Spearman correlation. For Pearson,
\begin{equation}
GB =
\frac{\sum_i(x_i-\bar x)(y_i-\bar y)}
{\sqrt{\sum_i(x_i-\bar x)^2}\sqrt{\sum_i(y_i-\bar y)^2}}.
\end{equation}

\section{Domain 5: Technical-supplement Fairness}

\subsection{Metric 15: Intersectional Calibration Error}

For group $g$ and bin $b$, let $S_{gb}$ be samples in intersectional group $g$
whose confidence falls in bin $b$. The implementation uses equal-width bins on
$[0,1]$. This extends expected calibration error as used in neural-network
calibration studies \cite{guo2017}. Let
\begin{align}
acc(S_{gb}) &= \frac{1}{|S_{gb}|}\sum_{i\in S_{gb}}Z_i,\\
conf(S_{gb}) &= \frac{1}{|S_{gb}|}\sum_{i\in S_{gb}}C_i.
\end{align}
The group-specific calibration error is
\begin{equation}
ECE_g = \sum_b \frac{|S_{gb}|}{|S_g|}
\left|acc(S_{gb})-conf(S_{gb})\right|.
\end{equation}
The intersectional calibration error is the population-weighted average
\begin{equation}
ICE = \sum_g \frac{|S_g|}{\sum_h |S_h|}ECE_g.
\end{equation}
The maximum calibration gap returned as \texttt{delta\_ice} is
\begin{equation}
\Delta ICE = \max_g ECE_g-\min_g ECE_g.
\end{equation}

\subsection{Metric 16: Weighted Clinical Harm-adjusted Fairness Gap}

Let $w_i$ be a clinical severity weight and $\ell_i=L(\hat Y_i,Y_i)$ be a
per-sample loss. For group $g$,
\begin{equation}
H(g)=\frac{1}{n_g}\sum_{i:G_i=g}w_i\ell_i.
\end{equation}
The implemented maximum gap is
\begin{equation}
wHAFG_{\max}=\max_g H(g)-\min_g H(g).
\end{equation}

\subsection{Metric 17: Counterfactual Parity Score}

Let $s_i\in[0,1]$ be a precomputed semantic similarity between the original
response and the response under a demographic swap. For a single pair type,
\begin{equation}
CPS=\frac{1}{n}\sum_{i=1}^{n}s_i.
\end{equation}
The demographic-swap construction is grounded in the counterfactual fairness
intuition that decisions should be stable under protected-attribute changes
when clinically relevant facts are unchanged \cite{kusner2017}.
For multiple swap-pair labels $p$, the library calculates
\begin{equation}
CPS_p=\frac{1}{n_p}\sum_{i\in p}s_i,\qquad
CPS=\frac{1}{\sum_p n_p}\sum_p\sum_{i\in p}s_i.
\end{equation}

\subsection{Metric 17a: Counterfactual Unfairness}

The implementation returns a complement to CPS. For a single pair,
\begin{equation}
CFU=1-CPS.
\end{equation}
For multiple pairs,
\begin{equation}
CFU=1-\min_p CPS_p.
\end{equation}

\subsection{Metric 18: Semantic Robustness Parity Index}

Let $R_i$ be a per-query robustness score, usually the mean similarity among
responses to semantically equivalent paraphrases. For group $g$,
\begin{equation}
R(g)=\frac{1}{n_g}\sum_{i:G_i=g}R_i.
\end{equation}
The implemented parity ratio is
\begin{equation}
SRPI=\frac{\min_g R(g)}{\max_g R(g)}.
\end{equation}
If the maximum group robustness is zero, the function returns zero.

\subsection{Metric 19: Lexical Diversity Disparity Index}

Let $T_g$ be the multiset of tokens pooled across all responses in group $g$ and
$V_g$ its vocabulary. The implemented root type-token ratio is
\begin{equation}
RTTR(g)=\frac{|V_g|}{\sqrt{|T_g|}}.
\end{equation}
The disparity index is
\begin{equation}
LDDI=\max_g RTTR(g)-\min_g RTTR(g).
\end{equation}
With $RTTR_{\mathrm{all}}$ computed after pooling all groups,
\begin{equation}
LDDI_{\mathrm{norm}}=\frac{LDDI}{RTTR_{\mathrm{all}}}.
\end{equation}

\subsection{Metric 20: Recommendation Entropy Gap}

Let $P_g(t)$ be the empirical distribution of recommendation labels $t$ in group
$g$. The group recommendation entropy is based on Shannon entropy
\cite{shannon1948}:
\begin{equation}
H(T\mid g)=-\sum_t P_g(t)\log_2 P_g(t).
\end{equation}
The implemented gap is
\begin{equation}
REG=\max_g H(T\mid g)-\min_g H(T\mid g).
\end{equation}
The implementation also returns
\begin{equation}
REG_{KL}=\max_g \sum_t P_g(t)\log_2\frac{P_g(t)}{P(t)},
\end{equation}
where $P(t)$ is the marginal recommendation distribution.

\subsection{Metric 21: Clinical Information Density Ratio}

For response $i$, let $a_i$ be the number of extracted clinical concepts and
$b_i$ the number of tokens. The response-level clinical information density is
\begin{equation}
CID_i=100\frac{a_i}{b_i}.
\end{equation}
For group $g$,
\begin{equation}
CID(g)=\frac{1}{n_g}\sum_{i:G_i=g}CID_i.
\end{equation}
The group ratio is
\begin{equation}
CIDR(g)=\frac{CID(g)}{\max_h CID(h)},
\end{equation}
and the reported parity summary is
\begin{equation}
CIDR_{\min}=\min_g CIDR(g).
\end{equation}
If the maximum group density is zero, all implemented ratios are set to zero.

\subsection{Metric 22: Diagnostic Completeness Index}

Let $D^\star$ be the reference differential diagnosis set. For response $i$,
\begin{equation}
DCI_i=\frac{|D(R_i)\cap D^\star|}{|D^\star|}.
\end{equation}
The group mean is
\begin{equation}
DCI(g)=\frac{1}{n_g}\sum_{i:G_i=g}DCI_i,
\end{equation}
and the implemented disparity is
\begin{equation}
\Delta DCI=\max_g DCI(g)-\min_g DCI(g).
\end{equation}
When diagnosis weights $w_d$ are supplied, the weighted response score is
\begin{equation}
wDCI_i=
\frac{\sum_{d\in D(R_i)\cap D^\star}w_d}{\sum_{d\in D^\star}w_d}.
\end{equation}

\subsection{Metric 23: Uncertainty Quantification Gap}

Let $h_i$ be the number of hedging terms in response $R_i$ and $q_i$ the number
of sentences. The uncertainty density is
\begin{equation}
UD_i=\frac{h_i}{q_i}.
\end{equation}
If no sentence is detected, $UD_i=0$. For group $g$,
\begin{equation}
UD(g)=\frac{1}{n_g}\sum_{i:G_i=g}UD_i.
\end{equation}
The implemented uncertainty quantification gap is
\begin{equation}
UQG=\max_g UD(g)-\min_g UD(g).
\end{equation}

\subsection{Metric 24: Geographic Representation Bias Index}

Let $p_c(r)$ be the normalized corpus evidence share in region $r$ and $p_b(r)$
the normalized burden share. EquiMed-DSS implements directed KL divergence in
nats \cite{kullback1951}:
\begin{equation}
GRBI=\KL(P_c\Vert P_b)=\sum_{r:p_c(r)>0}p_c(r)\log\frac{p_c(r)}{p_b(r)}.
\end{equation}
If $p_c(r)>0$ and $p_b(r)=0$, the implementation raises an error because KL is
undefined. With optional high-income regions $\mathcal{H}$, the high-income
overrepresentation ratio is
\begin{equation}
HIC_{\mathrm{ratio}}=
\frac{\sum_{r\in\mathcal{H}}p_c(r)}
{\sum_{r\in\mathcal{H}}p_b(r)}.
\end{equation}

\subsection{Metric 25: Healthcare System Stratified Fairness}

For system stratum $s$, let
\begin{equation}
\Delta_s=\max_g \E[Y\mid G=g,S=s]-\min_g \E[Y\mid G=g,S=s].
\end{equation}
The implemented within-system fairness gap is
\begin{equation}
HSSF=\sum_s P(S=s)\Delta_s.
\end{equation}
The implementation returns $\Delta_{\mathrm{within}}=HSSF$. It also returns the
population-weighted between-system variance
\begin{equation}
\Delta_{\mathrm{between}}=
\sum_s P(S=s)\left(\E[Y\mid S=s]-\sum_{u}P(S=u)\E[Y\mid S=u]\right)^2.
\end{equation}

\subsection{Metric 26: Intersectional Shapley Fairness Value}

Let $\mathcal{A}=\{A_1,\ldots,A_m\}$ be protected attributes. The attribution
formula follows Shapley's cooperative-game value \cite{shapley1953}. For a
subset $S\subseteq\mathcal{A}$, the implemented characteristic function is
\begin{equation}
v(S)=
\max_{a\in\mathrm{dom}(S)}\E[Y\mid A_S=a]
-
\min_{a\in\mathrm{dom}(S)}\E[Y\mid A_S=a],
\end{equation}
with $v(\varnothing)=0$. Cells below the configured minimum cell size are
ignored. The Shapley attribution for attribute $A_j$ is
\begin{equation}
\phi_j=
\sum_{S\subseteq\mathcal{A}\setminus\{A_j\}}
\frac{|S|!(m-|S|-1)!}{m!}
\left[v(S\cup\{A_j\})-v(S)\right].
\end{equation}
The total disparity is $v(\mathcal{A})$. Pairwise interactions are
\begin{equation}
I(A_j,A_k)=v(\{A_j,A_k\})-v(\{A_j\})-v(\{A_k\}).
\end{equation}

\section{Geographic Module}

\subsection{Metric 27: Burden-evidence Mismatch Index}

Let $e(r)$ be the normalized evidence share and $b(r)$ the normalized burden
share over the union of supplied regions. EquiMed-DSS defines BEMI as total
variation distance \cite{gibbs2002}:
\begin{equation}
BEMI=\frac{1}{2}\sum_{r\in\mathcal{R}}|e(r)-b(r)|.
\end{equation}
The per-region mismatch and ratio returned by the implementation are
\begin{align}
M(r)&=e(r)-b(r),\\
\rho(r)&=\frac{e(r)}{b(r)}.
\end{align}
The ratio $\rho(r)$ is finite only when $b(r)>0$; the implementation records a
missing value when the burden share is zero.
The most underserved region is the region with the minimum $M(r)$.

\subsection{Metric 28: Geographic Concentration of Coverage}

Let $x_r\ge 0$ be regional evidence counts for $R$ regions and
$p_r=x_r/\sum_u x_u$. The raw categorical Gini is
\begin{equation}
G_{\mathrm{raw}}=
\frac{\sum_{r=1}^{R}\sum_{u=1}^{R}|x_r-x_u|}
{2R\sum_{r=1}^{R}x_r}.
\end{equation}
Because the maximum raw Gini for $R$ categories is $(R-1)/R$, the implementation
uses the corrected value
\begin{equation}
G^\star=\frac{R}{R-1}G_{\mathrm{raw}}.
\end{equation}
The normalized Shannon entropy follows Shannon's entropy definition
\cite{shannon1948}:
\begin{equation}
H_{\mathrm{norm}}=
-\frac{\sum_{r:p_r>0}p_r\log p_r}{\log R},
\end{equation}
and the concentration score is
\begin{equation}
C_{\mathrm{geo}}=1-H_{\mathrm{norm}}.
\end{equation}

\section{Appendix Metrics}

\subsection{Metric 29: Bootstrap Confidence Interval}

Given data $x_1,\ldots,x_n$ and a statistic $T(\cdot)$, EquiMed-DSS draws
$B$ bootstrap resamples $x^{\ast b}$ of size $n$ with replacement and computes
\begin{equation}
\theta_b^\ast=T(x^{\ast b}),\qquad b=1,\ldots,B.
\end{equation}
For significance level $\alpha$, the implemented percentile interval is
\begin{equation}
CI_{1-\alpha}=
\left[
Q_{\alpha/2}(\theta_1^\ast,\ldots,\theta_B^\ast),
Q_{1-\alpha/2}(\theta_1^\ast,\ldots,\theta_B^\ast)
\right].
\end{equation}
The observed statistic is $T(x_1,\ldots,x_n)$. The percentile construction
follows the nonparametric bootstrap framework of Efron and Tibshirani
\cite{efron1979,efron1993}.

\subsection{Metric 30: Statistical Power Analysis}

The canonical implementation delegates two-sample $t$-test power and sample-size
calculations to \texttt{statsmodels.stats.power.tt\_ind\_solve\_power}. The
underlying planning target is Cohen's standardized effect
\begin{equation}
d=\frac{\mu_1-\mu_2}{\sigma}.
\end{equation}
The solver returns the per-group sample size $n$ satisfying
\begin{equation}
\mathrm{Power}=
P\left(\mathrm{reject}\ H_0:\mu_1=\mu_2\mid d,\alpha,n\right)
\end{equation}
for the requested alternative. The returned per-group sample size is
$\lceil n\rceil$, while the returned total sample size follows the implementation
as $\lceil 2n\rceil$. The standardized effect-size scale follows Cohen's
two-sample convention \cite{cohen1988}.

\subsection{Metric 31: Bias Concentration Index}

For nonnegative group bias proportions $p_1,\ldots,p_K$, the implementation
defines
\begin{equation}
BCI_{\mathrm{bias}}=
1-\frac{\sum_{g=1}^{K}p_g^2}{\left(\sum_{g=1}^{K}p_g\right)^2}.
\end{equation}
If the vector is empty or sums to zero, the value is zero. For a normalized
nonnegative vector, the finite-group upper bound is $1-1/K$, attained by an even
distribution. Larger values indicate more distributed bias; smaller values
indicate concentration in fewer groups.

\subsection{Metric 32: Mutual Information Content}

Let $D$ be a demographic categorical variable and $O$ an outcome categorical
variable. In the implementation, $D$ is expected to be encoded as nonnegative
integer categories for the entropy normalization. The raw mutual information
implemented via \texttt{mutual\_info\_score}
uses Shannon mutual information \cite{shannon1948}:
\begin{equation}
MIC=I(D;O)=\sum_{d,o}p(d,o)\log\frac{p(d,o)}{p(d)p(o)}.
\end{equation}
The implementation also reports a normalized value using demographic entropy:
\begin{equation}
MIC_{\mathrm{norm}}=\frac{I(D;O)}{H(D)},\qquad
H(D)=-\sum_d p(d)\log p(d).
\end{equation}
If $H(D)=0$, the normalized value is zero.

\subsection{Metric 33: Jensen-Shannon Divergence}

For two nonnegative vectors normalized to probability distributions $p$ and $q$,
let
\begin{equation}
m=\frac{p+q}{2}.
\end{equation}
The implementation returns the base-2 Jensen-Shannon divergence
\cite{lin1991}:
\begin{equation}
JSD(p,q)=
\frac{1}{2}\KL_2(p\Vert m)+\frac{1}{2}\KL_2(q\Vert m),
\end{equation}
where $\KL_2$ uses $\log_2$. The SciPy function returns the distance
$\sqrt{JSD}$, so the implementation squares that value. Both \texttt{jsd} and
\texttt{jsd\_distance} are reported.

\subsection{Metric 34: Wasserstein Distance}

For one-dimensional empirical samples $P_n=\{x_1,\ldots,x_n\}$ and
$Q_m=\{y_1,\ldots,y_m\}$, the implemented metric delegates to SciPy's first
Wasserstein distance without explicit sample weights
\cite{villani2009}:
\begin{equation}
WD(P_n,Q_m)=\inf_{\gamma\in\Gamma(P_n,Q_m)}
\int_{\R\times\R}|x-y|\,d\gamma(x,y),
\end{equation}
where $P_n$ and $Q_m$ are empirical distributions and $\Gamma(P_n,Q_m)$ is the
set of couplings with these marginals.
Equivalently, in one dimension,
\begin{equation}
WD(P_n,Q_m)=\int_{-\infty}^{\infty}|F_{P_n}(t)-F_{Q_m}(t)|\,dt.
\end{equation}

\subsection{Metric 35: Network Modularity}

Given an adjacency or correlation matrix $A$, the implementation constructs an
undirected graph using $|A|$ and detects communities with greedy modularity.
For total edge weight $2m=\sum_{ij}A_{ij}$, degree $k_i=\sum_j A_{ij}$, and
community assignment $c_i$, Newman modularity is \cite{newman2006,clauset2004}
\begin{equation}
Q=\frac{1}{2m}\sum_{i,j}
\left(A_{ij}-\frac{k_i k_j}{2m}\right)\ind(c_i=c_j).
\end{equation}

\subsection{Metric 36: Transparency Score}

For each explained decision $i$, the canonical appendix implementation expects
three scores in $[0,1]$: explanation quality $e_i$, feature importance $f_i$,
and interpretability $u_i$. This score is an implementation-level aggregate for
post-hoc explanation adequacy, aligned with the clinical need to expose reasons
for model outputs rather than predictions alone \cite{ribeiro2016}. The
per-decision transparency contribution is
\begin{equation}
t_i=\frac{e_i+f_i+u_i}{3}.
\end{equation}
The transparency score is the empirical mean
\begin{equation}
TS=\frac{1}{n}\sum_{i=1}^{n}t_i.
\end{equation}
If no explanations are provided, the implementation returns zero.

\subsection{Metric 37: Robustness Certification Score}

Let $a_{ib}$ be the agreement indicator between the original prediction for
case $i$ and the prediction under perturbation batch $b$:
\begin{equation}
a_{ib}=\ind(\hat Y_i=\hat Y_{ib}^{\mathrm{pert}}).
\end{equation}
For each perturbation batch,
\begin{equation}
r_b=\frac{1}{n}\sum_{i=1}^{n}a_{ib}.
\end{equation}
The implemented robustness certification score is
\begin{equation}
RCS=\frac{1}{B}\sum_{b=1}^{B}r_b.
\end{equation}
The implementation also reports
\begin{align}
SD(RCS) &= \sqrt{\frac{1}{B}\sum_{b=1}^{B}(r_b-RCS)^2},\\
r_{\min} &= \min_b r_b,\qquad r_{\max}=\max_b r_b.
\end{align}
The input argument $\epsilon$ is recorded in the output but is not used in the
calculation itself.

\section{Implemented Statistical Estimands Included to Reach 39 Derivations}

\subsection{Metric 38: Hierarchical Variance Partitioning and MAIHDA-style VPC}

For a Gaussian mixed model with outcome $Y_{ij}$ for individual $i$ in group $j$,
EquiMed-DSS fits a random-intercept model using the standard multilevel
variance-partitioning framework \cite{raudenbush2002}:
\begin{equation}
Y_{ij}=\beta_0+X_{ij}^{\top}\beta+u_j+\varepsilon_{ij},
\qquad
u_j\sim N(0,\sigma_u^2),\quad
\varepsilon_{ij}\sim N(0,\sigma_e^2).
\end{equation}
The implemented intraclass correlation from the null model is
\begin{equation}
ICC_{\mathrm{HLM}}=
\frac{\sigma_u^2}{\sigma_u^2+\sigma_e^2}.
\end{equation}
The reported pseudo-$R^2$ is the proportional reduction in total variance from
the null model to the full model:
\begin{equation}
R^2_{\mathrm{pseudo}}=
1-\frac{\sigma_{u,\mathrm{full}}^2+\sigma_{e,\mathrm{full}}^2}
{\sigma_{u,\mathrm{null}}^2+\sigma_{e,\mathrm{null}}^2}.
\end{equation}
For a binary outcome analysed on the logistic latent scale, the package
documentation notes the MAIHDA-style variance partition coefficient
\begin{equation}
VPC_{\mathrm{logit}}=
\frac{\sigma_u^2}{\sigma_u^2+\pi^2/3}.
\end{equation}
If the mixed-model fit fails, the implementation falls back to an ANOVA-style
decomposition with $J$ groups:
\begin{align}
SS_B &= \sum_j n_j(\bar Y_j-\bar Y)^2,\\
SS_W &= \sum_i (Y_i-\bar Y_{g(i)})^2,\\
MS_B &= \frac{SS_B}{J-1},\qquad
MS_W = \frac{SS_W}{n-J}.
\end{align}
With $\bar n$ denoting the mean group size, the fallback ICC is bounded to
$[0,1]$:
\begin{equation}
ICC_{\mathrm{fallback}}=
\max\left\{0,\min\left[1,
\frac{MS_B-MS_W}{MS_B+(\bar n-1)MS_W}
\right]\right\}.
\end{equation}

\subsection{Metric 39: Causal Mediation and Proportion Mediated}

Let $X$ be the treatment or exposure, $M$ the mediator, $Y$ the outcome, and
$C$ optional covariates. The implemented product-of-coefficients mediation uses
linear regression models \cite{sobel1982,imai2010}:
\begin{align}
M &= \alpha_0+\alpha_1X+C^\top\alpha_C+\varepsilon_M,\\
Y &= \beta_0+\beta_1X+\beta_2M+C^\top\beta_C+\varepsilon_Y.
\end{align}
A separate total-effect model is
\begin{equation}
Y=\tau_0+\tau_1X+C^\top\tau_C+\varepsilon_T.
\end{equation}
The indirect, direct, and total effects are
\begin{align}
IE &= \alpha_1\beta_2,\\
DE &= \beta_1,\\
TE &= \tau_1.
\end{align}
The implemented proportion mediated is
\begin{equation}
PM=\frac{IE}{TE},
\end{equation}
with $PM=0$ if $|TE|\le 10^{-10}$. The indirect-effect confidence interval is a
percentile bootstrap over the product $\alpha_1^\ast\beta_2^\ast$.
The Sobel test implemented in the same class is
\begin{align}
SE_{\mathrm{Sobel}} &=
\sqrt{\alpha_1^2SE(\beta_2)^2+\beta_2^2SE(\alpha_1)^2},\\
z_{\mathrm{Sobel}} &= \frac{\alpha_1\beta_2}{SE_{\mathrm{Sobel}}},\\
p &= 2\left[1-\Phi(|z_{\mathrm{Sobel}}|)\right].
\end{align}

\begin{thebibliography}{99}

\bibitem{shrout1979}
Shrout PE, Fleiss JL. Intraclass correlations: uses in assessing rater
reliability. \textit{Psychological Bulletin}. 1979;86(2):420-428.
doi:10.1037/0033-2909.86.2.420.

\bibitem{blandaltman1986}
Bland JM, Altman DG. Statistical methods for assessing agreement between two
methods of clinical measurement. \textit{Lancet}. 1986;327(8476):307-310.
doi:10.1016/S0140-6736(86)90837-8.

\bibitem{wilson1927}
Wilson EB. Probable inference, the law of succession, and statistical
inference. \textit{Journal of the American Statistical Association}.
1927;22(158):209-212. doi:10.1080/01621459.1927.10502953.

\bibitem{uniformguidelines1978}
Equal Employment Opportunity Commission, Civil Service Commission, Department
of Labor, Department of Justice. Uniform Guidelines on Employee Selection
Procedures. \textit{Federal Register}. 1978;43(166):38290-38315.

\bibitem{gini1912}
Gini C. \textit{Variabilita e Mutabilita}. Bologna: Tipografia di Paolo
Cuppini; 1912.

\bibitem{shewhart1931}
Shewhart WA. \textit{Economic Control of Quality of Manufactured Product}.
New York: D. Van Nostrand Company; 1931.

\bibitem{guo2017}
Guo C, Pleiss G, Sun Y, Weinberger KQ. On calibration of modern neural
networks. In: \textit{Proceedings of the 34th International Conference on
Machine Learning}. PMLR. 2017;70:1321-1330.

\bibitem{kusner2017}
Kusner MJ, Loftus J, Russell C, Silva R. Counterfactual fairness. In:
\textit{Advances in Neural Information Processing Systems}. 2017;30.

\bibitem{shannon1948}
Shannon CE. A mathematical theory of communication. \textit{The Bell System
Technical Journal}. 1948;27(3):379-423, 27(4):623-656.

\bibitem{kullback1951}
Kullback S, Leibler RA. On information and sufficiency. \textit{Annals of
Mathematical Statistics}. 1951;22(1):79-86.
doi:10.1214/aoms/1177729694.

\bibitem{lin1991}
Lin J. Divergence measures based on the Shannon entropy. \textit{IEEE
Transactions on Information Theory}. 1991;37(1):145-151.
doi:10.1109/18.61115.

\bibitem{gibbs2002}
Gibbs AL, Su FE. On choosing and bounding probability metrics.
\textit{International Statistical Review}. 2002;70(3):419-435.
doi:10.1111/j.1751-5823.2002.tb00178.x.

\bibitem{efron1979}
Efron B. Bootstrap methods: another look at the jackknife. \textit{Annals of
Statistics}. 1979;7(1):1-26. doi:10.1214/aos/1176344552.

\bibitem{efron1993}
Efron B, Tibshirani RJ. \textit{An Introduction to the Bootstrap}. New York:
Chapman and Hall/CRC; 1993.

\bibitem{cohen1988}
Cohen J. \textit{Statistical Power Analysis for the Behavioral Sciences}. 2nd
ed. Hillsdale: Lawrence Erlbaum Associates; 1988.

\bibitem{villani2009}
Villani C. \textit{Optimal Transport: Old and New}. Berlin: Springer; 2009.

\bibitem{newman2006}
Newman MEJ. Modularity and community structure in networks. \textit{Proceedings
of the National Academy of Sciences of the United States of America}.
2006;103(23):8577-8582. doi:10.1073/pnas.0601602103.

\bibitem{clauset2004}
Clauset A, Newman MEJ, Moore C. Finding community structure in very large
networks. \textit{Physical Review E}. 2004;70:066111.
doi:10.1103/PhysRevE.70.066111.

\bibitem{shapley1953}
Shapley LS. A value for n-person games. In: Kuhn HW, Tucker AW, eds.
\textit{Contributions to the Theory of Games II}. Princeton: Princeton
University Press; 1953:307-317.

\bibitem{ribeiro2016}
Ribeiro MT, Singh S, Guestrin C. ``Why should I trust you?'': explaining the
predictions of any classifier. In: \textit{Proceedings of the 22nd ACM SIGKDD
International Conference on Knowledge Discovery and Data Mining}. 2016:1135-1144.
doi:10.1145/2939672.2939778.

\bibitem{raudenbush2002}
Raudenbush SW, Bryk AS. \textit{Hierarchical Linear Models: Applications and
Data Analysis Methods}. 2nd ed. Thousand Oaks: Sage Publications; 2002.

\bibitem{sobel1982}
Sobel ME. Asymptotic confidence intervals for indirect effects in structural
equation models. \textit{Sociological Methodology}. 1982;13:290-312.
doi:10.2307/270723.

\bibitem{imai2010}
Imai K, Keele L, Tingley D. A general approach to causal mediation analysis.
\textit{Psychological Methods}. 2010;15(4):309-334. doi:10.1037/a0020761.

\end{thebibliography}

\section*{Implementation Notes}

\begin{itemize}[leftmargin=*]
\item BEMI is total variation distance, not one minus a correlation.
\item GRBI is directed KL divergence in nats, not a symmetric distance.
\item JSD is reported as divergence in base 2. The corresponding distance is
also exposed as the square root.
\item ECS is implemented as cosine distance, so lower values mean greater
consistency.
\item HAFG in \texttt{domain2} is a two-group count-based metric normalized by
the larger harm. wHAFG in \texttt{domain5} is a per-sample,
severity-weighted multi-group generalization.
\item Bootstrap confidence intervals in the canonical appendix implementation
use percentile intervals only.
\item The prior PDF derivations for SPG, CHR, IVI, GRI, ICE, wHAFG, LDDI, and
related metrics are broadly consistent with the implementation, but this file
uses the exact implemented details where the older derivations were more
general.
\end{itemize}

\end{document}
